% ID: SEL-D0073-OA-E03
% Tanda S073.
\textbf{Parte I. a)} Escribimos primero la nacionalidad de Laura y luego la de Vicente. Con $F$, $I$ y $E$ para francesa, italiana y española, el espacio es
\[
 \Omega=\{(F,F),(F,I),(I,F),(I,I),(E,F),(E,I)\}.
\]
Los pares no son equiprobables. Como cada persona extrae de su propio monedero al azar, multiplicamos las probabilidades de las extracciones independientes.

\textbf{b)} Usamos el complementario de que coincidan en nacionalidad:
\[
 1-\left(\frac{6}{12}\frac{9}{12}+\frac{2}{12}\frac{3}{12}\right)
 =\boxed{\frac{7}{12}}.
\]
\textbf{c)} Laura tiene seis monedas no francesas y Vicente tiene tres. La probabilidad es $(6/12)(3/12)=\boxed{1/8}$.

\textbf{Parte II. a)} La proporción muestral es $\hat p=45/300=0{,}15$. Hay 45 éxitos y 255 fracasos, por lo que usamos la aproximación normal. Para el $97\%$ de confianza, queda $0{,}015$ en cada cola y $z_{0{,}985}\approx2{,}1701$:
\[
 0{,}15\pm2{,}1701\sqrt{\frac{0{,}15\cdot0{,}85}{300}}
 \approx\boxed{[0{,}1053,\ 0{,}1947]}.
\]
\textbf{b)} El margen disminuye al $95\%$ porque el valor crítico baja a $1{,}96$ y el error estándar muestral permanece igual.
